\documentclass[10pt,twocolumn]{article}

\usepackage[margin=0.75in]{geometry}
\usepackage{amsmath}
\usepackage{amssymb}
\usepackage{booktabs}
\usepackage{array}
\usepackage{tabularx}
\usepackage{times}
\usepackage{url}

\pagestyle{plain}
\providecommand{\keywords}[1]{\par\smallskip\noindent\textbf{Keywords: }#1\par}
\newcolumntype{Y}{>{\raggedright\arraybackslash}X}

\title{Autonomous Construction of Ontology Function Layers from Operational Artifacts}

\author{Anonymous Authors\\Anonymous Institution}

\begin{document}
\maketitle

\begin{abstract}
Operational ontologies increasingly need an executable function layer in addition to object types, properties, relations, and constraints. This layer specifies which atomic operations can be invoked over ontology objects, what parameters they require, under which states they are valid, and what effects they produce. Existing automatic ontology construction methods mainly induce descriptive schemas and triples, while tool-extraction methods mainly produce callable API schemas. Neither directly constructs ontology-bound functions whose semantics are automatically validated without per-function expert approval. We formulate autonomous ontology Function Layer construction: given operational artifacts such as OpenAPI files, API documentation, SDK signatures, request-response examples, execution traces, and an object-layer draft, construct a versioned library of atomic functions grounded in ontology objects and properties. We propose an evidence-triangulated construction framework that discovers function candidates, jointly aligns functions to object-layer elements, generates evidence-bearing contracts, validates them through ontology constraints and replay evidence, and publishes or abstains according to calibrated risk. The evaluation measures semantic correctness, executability, selective publication risk, robustness to object-layer noise, and downstream workflow composability.
\end{abstract}

\keywords{Ontology construction, function layer, API semantics, autonomous validation, operational ontology}

\section{Introduction}
Enterprise ontologies are no longer only descriptive data models. In operational settings they must also define what computations, queries, and state-changing operations can be executed over domain objects. An object ontology may represent \textit{Order}, \textit{Customer}, and \textit{PaymentMethod}; an operational ontology additionally needs functions that retrieve an order, cancel a pending order, calculate a refund, or update a delivery address under defined state and policy constraints. Industrial ontology platforms expose similar function concepts for object-centric computation and actions, but constructing these functions remains a substantial engineering task.

Current research only partially addresses this requirement. Automatic ontology and knowledge-graph construction methods generate schemas, triples, and ontology drafts from text or data~\cite{kommineni2024human2machine,bai2025autoschemakg,lippolis2025ontogenia}. API-oriented methods extract tool schemas or API knowledge graphs from documentation~\cite{ni2025toolfactory,sun2025apikg,lee2025innout}. These lines of work remain separated. Ontology construction usually stops before executable operations, while API extraction usually stops at syntax-level callability. A JSON function schema can tell an agent that an API takes an \texttt{order\_id}, but it does not necessarily specify that the function operates on an \textit{Order}, requires a particular state, returns properties compatible with the ontology, or causes a known state transition.

This paper treats the Function Layer as a missing stage of automatic operational ontology construction. The input is not a single API page. It is a collection of fallible operational artifacts: database and object metadata, OpenAPI descriptions, prose API documentation, SDK signatures, request-response examples, error specifications, and observed executions. The output is not merely a tool schema, but an \textit{OntologyFunction} record with object bindings, parameter roles, property links, preconditions, effects, endpoint bindings, error contracts, evidence, confidence, and publication status.

The central constraint is no per-function human intervention during construction. Human labels are still required for evaluation, but the construction algorithm cannot ask an expert to accept or repair every candidate. This makes the problem scientific rather than only engineering: if human adjudication is removed, the system must decide which semantic claims are sufficiently supported by independent evidence and which ones should be rejected or abstained.

We propose an evidence-triangulated autonomous construction framework. Candidate functions are discovered from operational artifacts, aligned jointly to object types and properties, converted into evidence-bearing contracts, validated through ontology and cross-artifact checks, optionally checked against execution traces or replay harnesses, and then published using a calibrated selective-risk policy. The method can propose local object-layer corrections, but it does not silently rewrite the ontology without independent evidence.

The contributions are:
\begin{itemize}
    \item We define autonomous Function Layer construction as an operational ontology task distinct from descriptive ontology learning and API tool-schema extraction.
    \item We propose a joint object--function alignment and evidence-triangulated validation framework for ontology-bound atomic functions.
    \item We introduce an autonomous publication protocol that reports coverage--risk behavior instead of assuming every generated function is safe to publish.
    \item We specify an evaluation design covering field-level semantics, executability, calibration, object-layer noise, and downstream workflow composability.
\end{itemize}

\section{Related Work}
\subsection{Automatic Ontology and Knowledge-Graph Construction}
Recent LLM-supported ontology construction systems generate ontology drafts from requirements, induce schemas from text, or construct ontology-grounded knowledge graphs~\cite{kommineni2024human2machine,lippolis2025ontogenia,bai2025autoschemakg,feng2024ontologykg}. Text2KGBench evaluates fact generation under ontology constraints~\cite{mihindukulasooriya2023text2kg}. These works motivate automatic construction but primarily target descriptive knowledge. They do not evaluate whether operational functions are correctly bound to object types, properties, states, and executable endpoints.

\subsection{API Knowledge Graphs and Tool Extraction}
API knowledge graphs model endpoints, parameters, dependencies, and invocation relations for tool use and composition. Explore--Construct--Filter automatically constructs API knowledge graphs from documentation~\cite{sun2025apikg}, and In-N-Out models parameter-level API dependencies~\cite{lee2025innout}. ToolFactory extracts AI-compatible tools from REST API documentation and provides an API extraction benchmark~\cite{ni2025toolfactory}. ToolLLM and Gorilla study large-scale API selection and invocation~\cite{qin2024toolllm,patil2024gorilla}. These works are close to our setting, but the target is not an ontology Function Layer with object semantics, state preconditions, effects, and selective publication risk.

\subsection{Ontology-to-API and Semantic Service Composition}
Ontology-based API generation and semantic service composition compile ontology semantics into executable interfaces or compose APIs through typed input-output relations~\cite{garijo2020oba,callaghan2023biothings}. This direction is complementary. It assumes the ontology semantics are authoritative and generates or composes services from them. We solve the reverse problem: infer the operational function layer from existing artifacts whose semantics may be incomplete, inconsistent, or stale.

\section{Method}
\subsection{Problem Definition}
Let $A=\{A_1,\ldots,A_m\}$ denote operational artifacts and $O_0$ an object-layer draft. The goal is to construct a set of ontology functions:
\[
F = \{f_1,\ldots,f_n\}
\]
where each function is represented as:
\[
f=\langle id,T,k,I,O,P,E,B,C,\Pi,\gamma,\sigma\rangle.
\]
Here $T$ is the bound object type set, $k$ is the operation kind, $I$ and $O$ are typed inputs and outputs, $P$ are preconditions, $E$ are effects, $B$ is the endpoint binding, $C$ is the error contract, $\Pi$ is evidence, $\gamma$ is confidence, and $\sigma$ is publication status.

The algorithm must output published, rejected, and abstained candidates without expert decisions in the loop. The evaluation uses hidden gold labels but the method does not.

\subsection{Multi-Source Function Discovery}
Parsers extract operation candidates from OpenAPI files, API prose, SDK signatures, database procedures, examples, logs, and traces. For each candidate the system records endpoint identifiers, request and response fields, documented errors, source locations, and version metadata. Weak textual evidence is retained but assigned lower initial reliability.

\subsection{Joint Object--Function Alignment}
Freezing the object layer before function extraction causes every object-layer error to propagate. We therefore jointly score function, object, property, and parameter-role assignments. The score combines lexical similarity, dense similarity, datatype compatibility, request-response structure, graph neighborhood, observed read/write sets, and cross-function consistency. The output can bind to an existing ontology element, propose an alias or missing property, or mark the field unresolved.

\subsection{Evidence-Bearing Contract Generation}
The generator produces contracts only from local evidence and ontology neighborhoods. Critical fields such as target object, endpoint binding, required inputs, preconditions, and external effects must carry source references. If evidence is insufficient, the field remains unresolved rather than being completed from model prior knowledge.

\subsection{Automatic Validation}
Validation has three layers. Ontology validation checks domain, range, datatype, relation path, and state-predicate consistency. Cross-artifact validation checks whether prose, OpenAPI, examples, errors, and traces agree on required inputs, outputs, and effects. Execution validation invokes a sandbox or replay harness when available and checks request validity, response projection, error behavior, and observed state transitions.

\subsection{Selective Publication}
A publication score combines evidence diversity, ontology consistency, cross-artifact agreement, execution support, unresolved critical fields, and model disagreement. Thresholds are calibrated on a development set. Candidates above the publish threshold enter the Function Layer, candidates below the reject threshold are rejected, and intermediate candidates are abstained. The primary reported object is a risk--coverage curve, not a single cherry-picked threshold.

\section{Experiments}
\subsection{Datasets}
The experiment uses multiple resources because no public dataset fully annotates ontology-bound functions. ToolFactory provides heterogeneous API documentation and extraction labels~\cite{ni2025toolfactory}. Public OpenAPI repositories such as APIs.guru provide scalable endpoint and schema variation. API KG resources and parameter-dependency data support dependency evaluation~\cite{sun2025apikg,lee2025innout}. APIGen, BFCL, and function-calling datasets provide additional callability and argument-format checks. We add a Function Layer overlay with object bindings, parameter roles, preconditions, effects, error contracts, and test cases.

\subsection{Object-Layer Conditions}
We evaluate three object-layer settings: a gold object layer, an automatically generated object layer, and a controlled noisy object layer. Noise includes renamed object types, missing properties, wrong domain or range, merged classes, split classes, and incorrect relation edges. This is necessary because Paper 1 claims autonomous ontology construction, not only API schema extraction under a perfect ontology.

\subsection{Baselines}
Baselines include rule-based OpenAPI conversion, direct LLM document-to-tool extraction, ToolFactory-style extraction, API KG construction without ontology integration, sequential object-then-function alignment, our method without execution validation, our method without joint object correction, our method without abstention, and the full method. Human ontology engineering is reported as a cost reference, not as an autonomous baseline.

\subsection{Metrics}
We report function discovery precision, recall, and F1; endpoint exact match; target-object accuracy; property-link accuracy; parameter-role F1; precondition F1; effect F1; error-contract F1; full-function exact match; valid invocation rate; response-to-object projection accuracy; state-transition consistency; published coverage; published error rate; risk--coverage area; expected calibration error; unsupported-claim rate; degradation slope under object-layer noise; object-error amplification ratio; and downstream workflow executable rate.

\subsection{Analysis}
The main claim is supported only if the method publishes non-trivial coverage at a bounded error rate and improves ontology semantics and executability over tool-schema baselines. High field F1 with very low published coverage is not sufficient. The analysis traces errors across artifact parsing, alignment, contract generation, validation, publication, and downstream workflow composition.

\section{Conclusion}
This paper reframes atomic function extraction as autonomous operational ontology construction. The contribution is not another API schema extractor, but a no-review process that constructs ontology-bound functions, validates their semantics through independent evidence, and exposes a measurable coverage--risk boundary.

\begin{thebibliography}{99}

\bibitem{kommineni2024human2machine}
V. K. Kommineni et al.
\newblock From Human Experts to Machines: An LLM Supported Approach to Ontology and Knowledge Graph Construction.
\newblock arXiv:2403.08345, 2024.

\bibitem{bai2025autoschemakg}
Y. Bai et al.
\newblock AutoSchemaKG: Autonomous Schema and Knowledge Graph Construction.
\newblock arXiv:2505.23628, 2025.

\bibitem{lippolis2025ontogenia}
A. Lippolis et al.
\newblock OntoGenia: Ontology Generation with Large Language Models.
\newblock arXiv:2503.05388, 2025.

\bibitem{feng2024ontologykg}
S. Feng et al.
\newblock Ontology-Guided Knowledge Graph Construction with Large Language Models.
\newblock arXiv, 2024.

\bibitem{mihindukulasooriya2023text2kg}
N. Mihindukulasooriya et al.
\newblock Text2KGBench: A Benchmark for Ontology-Driven Knowledge Graph Generation.
\newblock arXiv:2308.02357, 2023.

\bibitem{ni2025toolfactory}
J. Ni et al.
\newblock ToolFactory: Automated Tool Extraction from API Documentation.
\newblock arXiv:2501.16945, 2025.

\bibitem{sun2025apikg}
R. Sun et al.
\newblock Explore--Construct--Filter: Automated API Knowledge Graph Construction.
\newblock arXiv:2502.13412, 2025.

\bibitem{lee2025innout}
J. Lee et al.
\newblock In-N-Out: API Graphs for Parameter-Level Tool Dependencies.
\newblock arXiv:2509.01560, 2025.

\bibitem{qin2024toolllm}
Y. Qin et al.
\newblock ToolLLM: Facilitating Large Language Models to Master Real-world APIs.
\newblock In \emph{ICLR}, 2024.

\bibitem{patil2024gorilla}
S. Patil et al.
\newblock Gorilla: Large Language Model Connected with Massive APIs.
\newblock In \emph{NeurIPS}, 2024.

\bibitem{garijo2020oba}
D. Garijo and M. Osorio.
\newblock OBA: An Ontology-Based Framework for Creating REST APIs.
\newblock arXiv:2007.09206, 2020.

\bibitem{callaghan2023biothings}
B. Callaghan et al.
\newblock BioThings Explorer: A Query Engine for a Federated API Network.
\newblock arXiv:2304.09344, 2023.

\end{thebibliography}

\end{document}
